\section{Giới thiệu}
\subsection{Đặt vấn đề}

Trong bối cảnh nền kinh tế số và tinh thần khởi nghiệp (startup) đang bùng nổ mạnh mẽ, nhu cầu về một môi trường làm việc linh hoạt, chuyên nghiệp và giàu tính kết nối ngày càng trở nên cấp thiết. Trước đây tại các quán cafe chưa áp dụng mô hình coworking space thường xuyên gặp vấn đề khi kinh doanh đối với những khách hàng ngồi hàng giờ với một cốc cafe duy nhất. Thậm chí, đã có những quán cafe đã phải buộc khách hàng phải order thêm sản phẩm, điều này dẫn tới sự không hài lòng cho khách hàng. Thấu hiểu được nhu cầu cần một không gian làm việc tại quán cafe của thị trường, mô hình coworking space cafe ra đời và đã được khách hàng cực kỳ quan tâm và hưởng ứng. Sự chuyển dịch từ mô hình văn phòng truyền thống sang các không gian làm việc chung (Co-working Space) không chỉ là xu hướng nhất thời mà đã trở thành một nhu cầu thực tế của các freelancer, các doanh nghiệp nhỏ và cả những người làm việc từ xa. Tuy nhiên, nhiều người dùng hiện nay vẫn đang đứng trước sự lựa chọn khó khăn giữa việc làm việc tại các quán cà phê hay sử dụng dịch vụ văn phòng chia sẻ.

Thực tế cho thấy, dù các quán cà phê mang lại sự tự do và chi phí thấp, nhưng chúng thường tồn tại nhiều hạn chế về tiếng ồn, tính riêng tư, sự thiếu hụt các tiện ích văn phòng chuyên dụng (phòng họp, in ấn) và đặc biệt là thiếu đi một môi trường thúc đẩy năng suất lao động. Ngược lại, các hệ thống quản lý Co-working Space hiện nay tại Việt Nam dù đã phát triển nhưng vẫn đối mặt với nhiều thách thức: quy trình đặt chỗ và thanh toán còn thủ công, dữ liệu tình trạng không gian chưa được cập nhật theo thời gian thực, và quan trọng nhất là sự thiếu hụt các công cụ giúp gắn kết cộng đồng thành viên – giá trị cốt lõi vốn tạo nên sự khác biệt cho mô hình này.

Nhiều hệ thống hiện tại mới chỉ dừng lại ở mức độ là một "công cụ cho thuê bàn ghế" điện tử, thiếu đi sự tương tác giữa các cá nhân trong cùng một không gian. Điều này khiến các thành viên bỏ lỡ cơ hội tìm kiếm đối tác, chia sẻ nguồn lực và xây dựng mạng lưới chuyên môn ngay tại nơi mình làm việc. Xuất phát từ những yêu cầu thực tiễn đó, nhóm thực hiện đề tài nghiên cứu, phân tích, thiết kế và triển khai Hệ thống quản lý văn phòng chia sẻ tích hợp kết nối cộng đồng. Đề tài kỳ vọng mang lại giải pháp quản trị tập trung, tích hợp thanh toán tự động, đồng thời kiến tạo một nền tảng số giúp tối ưu hóa việc sử dụng không gian và thu hẹp khoảng cách giữa các thành viên thông qua tính năng gợi ý đối tác thông minh.

\subsection{Các hướng giải quyết liên quan}

Để đáp ứng nhu cầu quản lý và vận hành không gian làm việc chung, hiện nay đã có nhiều mô hình và giải pháp được triển khai. Tuy nhiên, mỗi hướng tiếp cận đều tồn tại những ưu điểm và hạn chế riêng:

\subsubsection*{1. Sử dụng các quán cà phê làm việc (Work-friendly Cafe)}
Đây là lựa chọn phổ biến nhất của các cá nhân làm việc tự do nhờ tính sẵn có và chi phí thấp. Tuy nhiên, hướng giải quyết này không thể coi là một giải pháp bền vững cho doanh nghiệp hay dự án dài hạn vì:
\begin{itemize}
\item Không gian không ổn định, ô nhiễm tiếng ồn và thiếu tính chuyên nghiệp khi tiếp khách.
\item Cơ sở hạ tầng (Internet, điện, bàn ghế) không được thiết kế cho việc làm việc cường độ cao.
\item Thiếu hụt hoàn toàn các công cụ quản lý và khả năng kết nối mạng lưới chuyên môn có tổ chức.
\end{itemize}

\subsubsection*{2. Sử dụng các phần mềm quản lý văn phòng truyền thống (ERP/Office Management)}
Nhiều đơn vị lựa chọn các hệ thống quản trị doanh nghiệp lớn để áp dụng vào Co-working Space. Hướng tiếp cận này có ưu điểm về tính chặt chẽ trong quản lý nhân sự và tài chính, nhưng gặp thách thức:
\begin{itemize}
\item Quy trình quá cồng kềnh, không phù hợp với tính linh hoạt của đặt chỗ theo giờ/ngày.
\item Chi phí bản quyền và vận hành rất cao, yêu cầu đội ngũ kỹ thuật chuyên sâu để bảo trì.
\item Không hỗ trợ các tính năng tương tác cộng đồng dành cho khách hàng thuê bên ngoài.
\end{itemize}

\subsubsection*{3. Các nền tảng SaaS quản lý Co-working chuyên dụng quốc tế}
Các giải pháp như \textit{Nexudus}, \textit{OfficeRnD} hay \textit{Cobot} cung cấp đầy đủ các tính năng vận hành. Tuy nhiên, khi áp dụng tại thị trường Việt Nam:
\begin{itemize}
\item Khó khăn trong việc tích hợp các cổng thanh toán nội địa (MoMo, ZaloPay, VietQR).
\item Ngôn ngữ và quy trình hỗ trợ khách hàng chưa thực sự tối ưu cho người dùng Việt.
\item Chi phí đăng ký theo tháng tính trên số lượng user thường khá cao đối với các đơn vị vừa và nhỏ.
\end{itemize}

\subsubsection*{4. Xây dựng hệ thống quản lý tập trung hướng cộng đồng và tự động hóa}
Đây là hướng tiếp cận hiện đại mà đề tài lựa chọn, tập trung vào việc mô-đun hóa các chức năng để vừa đảm bảo vận hành, vừa tăng cường giá trị kết nối. Hệ thống được thiết kế bao gồm:
\begin{itemize}
\item Quản lý tài nguyên linh hoạt: Sơ đồ chỗ ngồi thời gian thực, quản lý đa chi nhánh.
\item Tự động hóa thanh toán và vận hành: Tích hợp các ví điện tử (MoMo/PayOS) và quy trình check-in/out bằng mã định danh.
\item Hệ sinh thái Networking: Module quản lý kỹ năng và gợi ý đối tác (Matching) dựa trên nhu cầu thực tế của thành viên.
\end{itemize}

Bên cạnh đó, một điểm then chốt trong giải pháp của đề tài là khả năng Gợi ý đối tác thông minh (Smart Partner Matching). Các hệ thống hiện nay thường chỉ cung cấp danh sách thành viên tĩnh, nhưng đề tài hướng tới:
\begin{itemize}
\item Cho phép người dùng tự định nghĩa "Hồ sơ năng lực" và "Nhu cầu hợp tác" thông qua hệ thống Tags chuyên sâu.
\item Thuật toán tìm kiếm và gợi ý dựa trên sự tương đồng về kỹ năng và vị trí địa lý (cùng chi nhánh).
\item Tích hợp bảng tin cộng đồng nội bộ để chia sẻ dự án và tìm kiếm nguồn lực nhanh chóng.
\end{itemize}

Tổng kết lại, các giải pháp hiện có dù đã đáp ứng được phần nào nhu cầu đặt chỗ cơ bản nhưng vẫn bỏ ngỏ tiềm năng kết nối con người – yếu tố giúp Co-working Space vượt xa mô hình quán cà phê hay văn phòng truyền thống. Vì vậy, đề tài lựa chọn hướng tiếp cận tích hợp toàn diện, lấy người dùng làm trung tâm để tạo ra một hệ thống quản lý văn phòng chia sẻ hiện đại và giàu tính thực tiễn.

\subsection{Mục tiêu và tính cấp thiết của đề tài}
Mục tiêu của đề tài là phát triển một hệ thống quản lý văn phòng chia sẻ (co‑working Space) nhằm khắc phục việc quản lý rời rạc và tận dụng lợi thế của dịch vụ trực tuyến để nâng cao hiệu suất vận hành và trải nghiệm khách hàng. Hệ thống bao gồm:
\begin{itemize}
\item Cho quản lý chuỗi: cổng web cho phép quản lý sản phẩm và tồn kho, quản lý chương trình khuyến mãi, đăng tuyển, quản lý chi nhánh, chỗ ngồi, phòng và thiết bị cho thuê, quản lý nhân sự, và thu thập/ứng xử với phản hồi khách hàng — tất cả được thiết kế để giảm sai sót thủ công, rút ngắn các quy trình vận hành và tối ưu hoá hoạt động tiếp thị.
\item Hỗ trợ đặt chỗ cho khách, ghi nhận và giao đơn đồ ăn/đồ uống, đặt/phát hành phòng và thiết bị cho thuê, theo dõi trạng thái đặt chỗ, xử lý hủy/đổi chỗ, và thanh toán tại chỗ — giúp nhân viên hoạt động nhanh, nhất quán và ít sai sót.
\item Cho phép tìm kiếm và đặt chỗ trước, xem chi tiết sản phẩm/dịch vụ, thuê phòng/thiết bị, theo dõi khuyến mãi và tuyển dụng, đánh giá dịch vụ, thanh toán trực tuyến an toàn, và tham gia diễn đàn cộng đồng để chia sẻ trải nghiệm và kết nối với người dùng khác.
\end{itemize}
Hệ thống hướng tới yêu cầu bảo mật, độ tin cậy và khả năng mở rộng để phục vụ nhiều chi nhánh, đồng thời tối ưu hóa quy trình nghiệp vụ từ đặt chỗ đến thanh toán và báo cáo.
\subsection{Phạm vi đề tài}
Sản phẩm đầu ra là một hệ thống quản lý văn phòng chia sẻ triển khai trong phạm vi Việt Nam, hoàn thiện trong vòng 2 kỳ theo kế hoạch phát triển. Tính năng cơ bản bao gồm cổng quản trị cho chủ/chủ quản lý chuỗi, cho nhân viên, và website dành cho khách hàng; các nền tảng hỗ trợ: trình duyệt phổ biến trên desktop và mobile. Đối tượng chính là nhóm 18–30 tuổi — những người thường tìm kiếm không gian làm việc linh hoạt, có khả năng chi trả cho dịch vụ và thành thạo công nghệ, phù hợp với mô hình co‑working space.